We compare SGD, SGD with weight decay (SGD+WD) and sharpness-aware minimisation (SAM) on a three-layer MLP trained on scikit-learn digits with symmetric label noise (0, 20, 40\%) and on a noisy two-spirals task, using five seeds, hyper-parameters selected on a noisy validation split with separate tuning seeds, and a fixed budget of 150 epochs (SAM therefore uses twice the gradient evaluations).
SAM clearly beats plain SGD under noise (clean test accuracy \rhval{main/sam/digits_n0.2/test_acc/mean:3} vs \rhval{main/sgd/digits_n0.2/test_acc/mean:3} at 20\% noise and \rhval{main/sam/digits_n0.4/test_acc/mean:3} vs \rhval{main/sgd/digits_n0.4/test_acc/mean:3} at 40\%), but most of this gain is also obtained by a tuned weight decay: the margin over SGD+WD is \rhval{cmp/main/sgd+wd/digits_n0.2/test_acc/delta:3} at 20\% noise (paired $p=\rhval{cmp/main/sgd+wd/digits_n0.2/test_acc/paired_p:3}$, Welch $p=\rhval{cmp/main/sgd+wd/digits_n0.2/test_acc/welch_p:2}$, so we call it inconclusive) and \rhval{cmp/main/sgd+wd/digits_n0.4/test_acc/delta:3} at 40\% (both $p<0.05$ uncorrected; one of five seeds is negative and the paired test does not survive multiple-comparison correction). On two-spirals SAM does not help.
On noisy digits accuracy rises with the radius $\rho$ up to $\rho=0.5$ and collapses to chance at $\rho=1.0$ (no value in between was tested); a random-direction perturbation of the same radius behaves like SGD.
A fixed-norm loss-increase sharpness is only weakly or inconsistently correlated with the generalisation gap (pooled Spearman \rhval{derived/corr-pooled-all/digits_n0.4/adv_loss/mean:2}--\rhval{derived/corr-pooled-all/digits_n0.4/rand_acc/mean:2} over the two sharpness measures and two gaps, negative on spirals); on noisy digits adversarial sharpness ranks SGD, which memorises the flipped labels, as flatter than the better-generalising SAM and SGD+WD, whereas random-direction sharpness ranks SAM flattest.
This is a small, single-architecture CPU study with reimplemented optimisers; we make no claim beyond this scale.
